\subsection{ADO'S THEOREM}

PROPOSITION 1. Let \(\mathfrak{g}\) be a Lie algebra, \(\mathfrak{n}\) its largest nilpotent ideal, \(\mathfrak{a}\) a nilpotent ideal of \(\mathfrak{g}\) and \(\rho\) a finite-dimensional representation of \(\mathfrak{a}\) such that every element of \(\rho(a)\) is nilpotent. Then \(\rho\) admits a finite-dimensional extension \(\sigma\) to \(\mathfrak{g}\) such that every element of \(\sigma(n)\) is nilpotent.

Let \(\mathfrak{a} = \mathfrak{n}_0 \subset \mathfrak{n}_1 \subset \cdots \subset \mathfrak{n}_p = \mathfrak{n}\) be a sequence of subalgebras of \(\mathfrak{n}\) such that \(\mathfrak{n}_{i-1}\) is an ideal of \(\mathfrak{n}_i\) of codimension 1 for \(1 \leqslant i \leqslant p\) (§ 4, no. 1, Proposition 1 (e)). The algebra \(\mathfrak{n}_i\) is therefore the direct sum of \(\mathfrak{n}_{i-1}\) and a 1-dimensional subalgebra. As \(\operatorname{ad}_x\) is nilpotent for all \(x \in \mathfrak{n}\) , it is possible (Theorem 1) to find one by one finite-dimensional extensions \(\rho_1, \rho_2, \ldots, \rho_p = \rho'\) of \(\rho\) to \(\mathfrak{n}_1, \mathfrak{n}_2, \ldots, \mathfrak{n}_p = \mathfrak{n}\) such that every element of \(\rho'(\mathfrak{n})\) is nilpotent.

Let \(\mathfrak{r}\) be the radical of \(\mathfrak{g}\) and let \(\mathfrak{n} = \mathfrak{r}_0 \subset \mathfrak{r}_1 \subset \dots \subset \mathfrak{r}_q = \mathfrak{r}\) be a sequence of subalgebras of \(\mathfrak{r}\) such that \(\mathfrak{r}_{i-1}\) is an ideal of \(\mathfrak{r}_i\) of codimension 1 for \(1 \leqslant i \leqslant q\) (§ 5, no. 1, Proposition 2 (d)). The algebra \(\mathfrak{r}_i\) is thus the direct sum of \(\mathfrak{r}_{i-1}\) and a 1-dimensional subalgebra. As \([\mathfrak{r}, \mathfrak{r}] \subset \mathfrak{n}\) , it is possible (Theorem 1) to find one by one finite-dimensional extensions \(\rho_1', \rho_2', \ldots, \rho_q' = \rho''\) of \(\rho'\) to \(\mathfrak{r}_1, \mathfrak{r}_2, \ldots, \mathfrak{r}_q = \mathfrak{r}\) such that every element of \(\rho''(\mathfrak{n})\) is nilpotent.

Finally g is the direct sum of r and a subalgebra s (§ 6, no. 8, Theorem 5). As \([s, r] \subset n\), it is possible (Theorem 1) to find a finite-dimensional extension σ of ρ'' to g such that every element of σ(n) is nilpotent.

THEOREM 2. Every Lie algebra has a finite-dimensional faithful linear representation.

More precisely:

THEOREM 3 (Ado). Let \(\mathfrak{g}\) be a Lie algebra and \(\mathfrak{n}\) its largest nilpotent ideal. There exists a finite-dimensional faithful representation \(\rho\) of \(\mathfrak{g}\) such that every element of \(\rho(\mathfrak{n})\) is nilpotent.

The 1-dimensional Lie algebra K admits finite-dimensional faithful representations \(\tau\) such that every element of \(\tau(\mathbf{K})\) is nilpotent, for example the representation

\[
\lambda \mapsto \left( \begin{array}{c c} 0 & 0 \\ \lambda & 0 \end{array} \right).
\]

It is easily deduced that the centre \(c\) of \(g\) , which is a product of 1-dimensional algebras, admits a finite-dimensional faithful representation \(\sigma\) such that every element of \(\sigma(c)\) is nilpotent. Let \(\sigma_1\) be a finite-dimensional extension of \(\sigma\) to \(g\) such that every element of \(\sigma_1(n)\) is nilpotent (Proposition 1); if \(t\) denotes the kernel of \(\sigma_1\) , then \(t \cap c = \{0\}\) . On the other hand, let \(\sigma_2\) be the adjoint representation of \(g\) , whose kernel is \(c\) ; every element of \(\sigma_2(n)\) is nilpotent. The direct sum \(\rho\) of \(\sigma_1\) and \(\sigma_2\) is finite-dimensional; every element of \(\rho(n)\) is nilpotent; and the kernel of \(\rho\) , contained in \(t\) and in \(c\) , is zero, so that \(\rho\) is faithful.

